\section{Experiment 1: Varying the Number of Expert Annotations}
\label{sec:experiment1}

\subsection{Overview}

Experiment 1 investigates how estimation quality changes as the number of expert-labeled samples $n$ grows, while the total dataset of $N$ LLM-annotated samples remains fixed. The central question is: given a budget of $n$ expert annotations, how much does each debiasing method improve over using only LLM labels, and how does performance compare to using only expert labels?

The experiment is structured in four phases of increasing model complexity, ranging from intercept-only prevalence estimation (Phase 1) to full multivariate logistic regression with five features (Phase 4). All phases share the same repetition structure: for each dataset--LLM pair and each value of $n$, the experiment is repeated $R = 300$ times with independently drawn expert subsets (seeds $1, \ldots, 300$). Metrics are then averaged over repetitions and standard errors reported.

\subsection{Datasets and LLMs}

Four datasets are used: \textsc{CUAD}, \textsc{Misogynistic}, \textsc{FOMC}, and \textsc{PubMedQA}. Five LLMs provide automated annotations: Llama, DeepSeek, GPT, Mistral, and Claude. Each dataset--LLM combination yields a fully annotated file of $N$ rows, where every item has both an LLM label $\hat{Y}_i$ and a ground-truth expert label $Y_i$. The expert labels are treated as the gold standard; the experiment simulates the scenario where only $n \ll N$ of these expert labels are available at estimation time.

\subsection{Ground Truth and Reference Estimator}

For each dataset--LLM pair, the ground truth parameter $\theta^*$ is computed using \emph{all} $N$ expert labels via unregularized logistic regression (or the logit of the sample mean for Phase~1). This serves as the oracle target that debiasing methods attempt to recover from the $n$-sample subset. The LLM-only baseline $\hat{\theta}_{\text{LLM}}$ is similarly computed from all $N$ LLM annotations.

\subsection{Expert Sampling}

At each repetition $r$ and sample size $n$, a random subset $\mathcal{L}_r \subset \{1, \ldots, N\}$ of size $n$ is drawn without replacement using seed $r$. The rows in $\mathcal{L}_r$ are treated as expert-labeled; the remaining $N - n$ rows are treated as unlabeled (LLM-only). The sample sizes $n$ are drawn from a log-spaced grid over $[20, 200]$, yielding approximately 10 values per run.

\subsection{Methods}

Four estimation methods are compared in every phase:

\paragraph{Expert-only ($\theta_\dagger$).}
A logistic regression model fitted exclusively on the $n$ expert-labeled rows using L2 regularization with $\lambda = 0.01$ (Phases 2--4) or the sample logit mean (Phase 1). This is the na\"{i}ve baseline that ignores all LLM annotations.

\paragraph{DSL.}
The Data Source Likelihood method \citep[see][]{dsl_paper}, fitted via the \texttt{dsl} R package. DSL models both expert and LLM labels jointly under a latent-variable framework, correcting for systematic LLM annotation error.

\paragraph{PPI.}
Prediction-Powered Inference \citep{angelopoulos2023prediction}. PPI uses the LLM annotations on all $N$ items to construct a bias-correction term, which is subtracted from the expert-only estimate. The correction is computed as the difference between the logistic loss on all LLM annotations and the logistic loss on the $n$ expert-labeled LLM annotations, then added to the expert loss:
\[
    \hat{\theta}_{\text{PPI}} = \arg\min_\theta \left[ (\mathcal{L}_{\text{unlab}}(\theta;\hat{Y}) - \mathcal{L}_{\text{lab}}(\theta;\hat{Y})) + \mathcal{L}_{\text{lab}}(\theta;Y) \right] + \frac{\lambda}{2}\|\beta_{1:p}\|^2.
\]

\paragraph{PPI++ (Adaptive PPI).}
PPI++ \citep{angelopoulos2023ppipp} extends PPI by introducing a scalar $\lambda^* \in [0,1]$ that adaptively weights the LLM correction term:
\[
    \hat{\theta}_{\text{PPI++}} = \arg\min_\theta \left[ \lambda^* (\mathcal{L}_{\text{unlab}}(\theta;\hat{Y}) - \mathcal{L}_{\text{lab}}(\theta;\hat{Y})) + \mathcal{L}_{\text{lab}}(\theta;Y) \right] + \frac{\lambda}{2}\|\beta_{1:p}\|^2.
\]
The weight $\lambda^*$ is estimated from the labeled data using the gradient covariance formula:
\[
    \lambda^* = \frac{\operatorname{tr}(V \Sigma_{\text{cov}} V)}{2(1 + n/N)\operatorname{tr}(V \Sigma_{\text{var}} V)},
    \quad \lambda^* \in [0, 1],
\]
where $V$ is the inverse of the regularized Hessian on unlabeled data, $\Sigma_{\text{cov}} = \frac{1}{n}\sum_{i \in \mathcal{L}} g_i g_{\hat{i}}^\top$ is the cross-covariance of expert and LLM gradients, and $\Sigma_{\text{var}} = \frac{1}{n}\sum_{i \in \mathcal{L}} g_{\hat{i}} g_{\hat{i}}^\top$ is the variance of LLM gradients. When the LLM is a reliable annotator, $\lambda^* \to 1$ and PPI++ reduces to standard PPI. When the LLM is unreliable, $\lambda^* \to 0$ and PPI++ reduces to expert-only.

For Phase 1 (intercept-only), PPI++ is applied with $\lambda_{\text{L2}} = 0$ and an empty feature matrix, so only the intercept $\beta_0$ is estimated.

\subsection{Phases}

\paragraph{Phase 1: Class Prevalence (Intercept-Only).}
No features are used. The target parameter is $\theta^* = \text{logit}(\bar{Y})$, the logit of the population-level class prevalence. All four methods estimate $\beta_0$ only.

\paragraph{Phase 2: Low-Variance Feature.}
A single feature $x_2$ with low across-sample variance is included. Each method estimates $\theta = [\beta_0, \beta_2]^\top$ using L2-regularized logistic regression ($\lambda = 0.01$). The feature is dataset-specific and defined in the dataset configuration.

\paragraph{Phase 3: High-Variance Feature.}
Identical to Phase 2 but using a single high-variance feature. This allows comparison of debiasing effectiveness across features with different signal strength.

\paragraph{Phase 4: Full Logistic Regression.}
All five features $x_1, \ldots, x_5$ are included. Each method estimates $\theta = [\beta_0, \beta_1, \ldots, \beta_5]^\top \in \mathbb{R}^6$ with L2 regularization ($\lambda = 0.01$).

\subsection{Evaluation Metrics}

Estimation quality is assessed using two metrics, computed in the space of probabilities (Phase 1) or coefficients (Phases 2--4).

\paragraph{Standardized RMSE (sRMSE).}
For Phase 1, estimation is performed in log-odds space but evaluation is in probability space to avoid division by zero on balanced datasets:
\[
    \text{sRMSE} = \sqrt{\frac{1}{R}\sum_{r=1}^{R} \left(\frac{\sigma(\hat{\beta}_0^{(r)}) - \sigma(\beta_0^*)}{\sigma(\beta_0^*)}\right)^2},
\]
where $\sigma(\cdot)$ is the sigmoid function.

For Phases 2--4, the Euclidean sRMSE over the full coefficient vector is used:
\[
    \text{sRMSE}_{\text{eucl}} = \sqrt{\frac{1}{R}\sum_{r=1}^{R} \frac{\|\hat{\theta}^{(r)} - \theta^*\|^2}{\|\theta^*\|^2}}.
\]

For Phases 2 and 3, an additional per-coefficient sRMSE is reported for the feature coefficient $\beta_2$:
\[
    \text{sRMSE}_{\beta_2} = \sqrt{\frac{1}{R}\sum_{r=1}^{R} \left(\frac{\hat{\beta}_2^{(r)} - \beta_2^*}{\beta_2^*}\right)^2}.
\]

\paragraph{Standardized Bias.}
\[
    \text{Bias} = \frac{1}{R}\sum_{r=1}^{R} \frac{\hat{\theta}^{(r)} - \theta^*}{\theta^*},
\]
computed analogously to sRMSE (in probability space for Phase 1, over $\beta_2$ for Phases 2--3, averaged over all coefficients for Phase 4).

Standard errors for both metrics are computed as $\text{SE} = \text{std}(\cdot) / \sqrt{R}$ over the $R = 300$ repetitions.

\subsection{Implementation Details}

All Python experiments are run via Apptainer containers on the Chalmers Vera HPC cluster using SLURM array jobs (array size 300, one seed per task). DSL is called via \texttt{rpy2} within each Python process, with the R session initialized once per job to avoid startup overhead. Results from each repetition are saved as \texttt{.npz} files and aggregated into CSV summary files by the evaluate scripts.

L2 regularization ($\lambda = 0.01$) is applied consistently across expert-only, PPI, and PPI++ in Phases 2--4. Phase 1 uses no regularization since the intercept-only model has no feature coefficients to penalize. The ground truth $\theta^*$ is computed with \emph{unregularized} logistic regression to avoid regularization bias in the reference target.
